\documentclass[11pt,a4paper]{article}
\newcommand{\doctitle}{Flood Probability Classifier}
\newcommand{\docsubtitle}{Gradient Boosting Machine Model Documentation}
\newcommand{\docversion}{2.2}
\newcommand{\docdate}{30-October-2026}
\newcommand{\docauthor}{David K Kelly}
% MKM Research Labs — Shared Document Header
% Usage: % MKM Research Labs — Shared Document Header
% Usage: % MKM Research Labs — Shared Document Header
% Usage: \input{../shared/mkm_header}  (before \begin{document})
%
% Required commands before \input:
%   \newcommand{\doctitle}{...}
%   \newcommand{\docsubtitle}{...}
%   \newcommand{\docversion}{1.0}
%   \newcommand{\docdate}{13-February-2026}
%   \newcommand{\docauthor}{David K Kelly}

% ─── Packages ───────────────────────────────────────────────────────────────
\usepackage[utf8]{inputenc}
\usepackage{amsmath,amsfonts,amssymb}
\usepackage{listings}
\usepackage{color,xcolor}
\usepackage{graphicx}
\usepackage{hyperref}
\usepackage[toc,page]{appendix}
\usepackage{tabto}
\usepackage{float}
\usepackage{comment}
\usepackage[en-GB]{datetime2}
\usepackage{setspace}
\usepackage{subcaption}
\usepackage{booktabs}
\usepackage{longtable}
\usepackage{adjustbox}
\usepackage{enumitem}
\usepackage{fancyhdr}
\usepackage{lastpage}

% ─── Page geometry ──────────────────────────────────────────────────────────
\onehalfspacing
\usepackage{geometry}
\geometry{a4paper, total={170mm,247mm}, left=20mm, top=25mm, bottom=25mm}

% ─── Headers and footers ───────────────────────────────────────────────────
\pagestyle{fancy}
\fancyhf{}
\fancyhead[L]{\small\textcolor{gray}{MKM Research Labs}}
\fancyhead[R]{\small\textcolor{gray}{\doctitle{} --- v\docversion}}
\fancyfoot[C]{\small\textcolor{gray}{CONFIDENTIAL}}
\fancyfoot[R]{\small\textcolor{gray}{Page \thepage\ of \pageref{LastPage}}}
\renewcommand{\headrulewidth}{0.4pt}
\renewcommand{\footrulewidth}{0.4pt}

% ─── Colours ────────────────────────────────────────────────────────────────
\definecolor{mkmblue}{HTML}{1565C0}
\definecolor{codegreen}{rgb}{0,0.6,0}
\definecolor{codegray}{rgb}{0.5,0.5,0.5}
\definecolor{codepurple}{rgb}{0.58,0,0.82}
\definecolor{backcolour}{rgb}{0.95,0.95,0.92}

% ─── Hyperref setup ─────────────────────────────────────────────────────────
\hypersetup{
    colorlinks=true,
    linkcolor=mkmblue,
    citecolor=mkmblue,
    urlcolor=mkmblue
}

% ─── Code listings ──────────────────────────────────────────────────────────
\lstdefinestyle{mystyle}{
    backgroundcolor=\color{backcolour},
    commentstyle=\color{codegreen},
    keywordstyle=\color{magenta},
    numberstyle=\tiny\color{codegray},
    stringstyle=\color{codepurple},
    basicstyle=\ttfamily\footnotesize,
    breakatwhitespace=false,
    breaklines=true,
    captionpos=b,
    keepspaces=true,
    numbers=left,
    numbersep=5pt,
    showspaces=false,
    showstringspaces=false,
    showtabs=false,
    tabsize=2
}
\lstset{style=mystyle}

% ─── Section formatting ────────────────────────────────────────────────────
\setcounter{tocdepth}{3}
\setcounter{secnumdepth}{3}

% ─── Shared macros ──────────────────────────────────────────────────────────
\newcommand{\mkmtitlepage}{%
    \begin{titlepage}
    \begin{center}
    \vspace*{1cm}
    \begin{figure}[H]\centering
        \includegraphics[width=0.4\textwidth]{../shared/Logos/MKM_logo.png}
    \end{figure}
    \vspace{1.5cm}
    {\Huge\bfseries\textcolor{mkmblue}{\doctitle}\par}
    \vspace{0.5cm}
    {\Large\textcolor{gray}{\docsubtitle}\par}
    \vspace{2cm}
    {\large MKM Research Labs\par}
    \vspace{0.5cm}
    {\large Version \docversion{} --- \docdate\par}
    \vspace{0.5cm}
    {\large \docauthor\par}
    \vspace{2cm}
    \textcolor{red}{\textbf{CONFIDENTIAL --- PROPRIETARY}}
    \end{center}
    \end{titlepage}
}

\newcommand{\mkmlegalpage}{%
    \clearpage
    \thispagestyle{empty}
    \begin{center}
    \large\textbf{Legal Notice}
    \vspace{1em}

    \textbf{PROPRIETARY AND CONFIDENTIAL}
    \end{center}

    \noindent This document, including all algorithms, models, methodology, formulas,
    calculations, and intellectual concepts contained herein, constitutes the exclusive
    intellectual property and confidential trade secrets of MKM Research Labs (``MKM'').

    \vspace{0.5em}
    \noindent Any usage, reproduction, distribution, modification, or disclosure of this
    document or any portion thereof, without the express prior written authorisation from
    MKM Research Labs is strictly prohibited and constitutes an infringement of intellectual
    property rights.

    \vspace{0.5em}
    \noindent The document is provided on an ``as is'' basis. MKM Research Labs makes no
    representations or warranties of any kind, express or implied, regarding its accuracy,
    reliability, or suitability for any particular purpose.

    \vspace{1em}
    \begin{center}
    All rights reserved. \copyright\ 2021--2026 MKM Research Labs.

    \vspace{0.5em}
    Governed by the laws of the United Kingdom.
    \end{center}
    \clearpage
}

\newcommand{\mkmhistorypage}[1]{%
    \section{Document History}
    \begin{table}[ht]
        \centering
        \begin{tabular}{l|l|l|l}
            \textbf{Date} & \textbf{Version} & \textbf{Description} & \textbf{Author} \\
            \hline
            #1
        \end{tabular}
        \label{tab:revision_history}
    \end{table}
    \clearpage
}

% ─── Bank validation protocol section commands ──────────────────────────────
% Use these to ensure every model doc has the required sections.
% They are just styled \section calls — content is provided by each doc.

\newcommand{\mkmexecsummary}{\section{Executive Summary}}
\newcommand{\mkmpurpose}{\section{Model Purpose and Scope}}
\newcommand{\mkmframework}{\section{Mathematical Framework}}
\newcommand{\mkminputs}{\section{Input Parameters and Calibration}}
\newcommand{\mkmimplementation}{\section{Implementation Details}}
\newcommand{\mkmsensitivity}{\section{Sensitivity Analysis}}
\newcommand{\mkmvalidation}{\section{Validation and Backtesting}}
\newcommand{\mkmtestresults}{\subsection{Automated Test Results}\input{test_results}}
\newcommand{\mkmlimitations}{\section{Model Limitations and Known Weaknesses}}

\newcommand{\mkmgovernance}{\section{Governance Validation Questions}
\noindent The following nine questions are mandated by Chapter~5 of the
\emph{Model Risk Governance for Vendors} handbook and must be explicitly
addressed for every model in the inventory. Response status is tracked in the
Model Governance Dashboard (MRC portal). \par\medskip}

\newcommand{\mkmvqitem}[3]{\subsection{Q#1 --- #2}
\noindent\textit{``#3''}\par\medskip\noindent}

\newcommand{\mkmchangelog}{\section{Change History}}
  (before \begin{document})
%
% Required commands before \input:
%   \newcommand{\doctitle}{...}
%   \newcommand{\docsubtitle}{...}
%   \newcommand{\docversion}{1.0}
%   \newcommand{\docdate}{13-February-2026}
%   \newcommand{\docauthor}{David K Kelly}

% ─── Packages ───────────────────────────────────────────────────────────────
\usepackage[utf8]{inputenc}
\usepackage{amsmath,amsfonts,amssymb}
\usepackage{listings}
\usepackage{color,xcolor}
\usepackage{graphicx}
\usepackage{hyperref}
\usepackage[toc,page]{appendix}
\usepackage{tabto}
\usepackage{float}
\usepackage{comment}
\usepackage[en-GB]{datetime2}
\usepackage{setspace}
\usepackage{subcaption}
\usepackage{booktabs}
\usepackage{longtable}
\usepackage{adjustbox}
\usepackage{enumitem}
\usepackage{fancyhdr}
\usepackage{lastpage}

% ─── Page geometry ──────────────────────────────────────────────────────────
\onehalfspacing
\usepackage{geometry}
\geometry{a4paper, total={170mm,247mm}, left=20mm, top=25mm, bottom=25mm}

% ─── Headers and footers ───────────────────────────────────────────────────
\pagestyle{fancy}
\fancyhf{}
\fancyhead[L]{\small\textcolor{gray}{MKM Research Labs}}
\fancyhead[R]{\small\textcolor{gray}{\doctitle{} --- v\docversion}}
\fancyfoot[C]{\small\textcolor{gray}{CONFIDENTIAL}}
\fancyfoot[R]{\small\textcolor{gray}{Page \thepage\ of \pageref{LastPage}}}
\renewcommand{\headrulewidth}{0.4pt}
\renewcommand{\footrulewidth}{0.4pt}

% ─── Colours ────────────────────────────────────────────────────────────────
\definecolor{mkmblue}{HTML}{1565C0}
\definecolor{codegreen}{rgb}{0,0.6,0}
\definecolor{codegray}{rgb}{0.5,0.5,0.5}
\definecolor{codepurple}{rgb}{0.58,0,0.82}
\definecolor{backcolour}{rgb}{0.95,0.95,0.92}

% ─── Hyperref setup ─────────────────────────────────────────────────────────
\hypersetup{
    colorlinks=true,
    linkcolor=mkmblue,
    citecolor=mkmblue,
    urlcolor=mkmblue
}

% ─── Code listings ──────────────────────────────────────────────────────────
\lstdefinestyle{mystyle}{
    backgroundcolor=\color{backcolour},
    commentstyle=\color{codegreen},
    keywordstyle=\color{magenta},
    numberstyle=\tiny\color{codegray},
    stringstyle=\color{codepurple},
    basicstyle=\ttfamily\footnotesize,
    breakatwhitespace=false,
    breaklines=true,
    captionpos=b,
    keepspaces=true,
    numbers=left,
    numbersep=5pt,
    showspaces=false,
    showstringspaces=false,
    showtabs=false,
    tabsize=2
}
\lstset{style=mystyle}

% ─── Section formatting ────────────────────────────────────────────────────
\setcounter{tocdepth}{3}
\setcounter{secnumdepth}{3}

% ─── Shared macros ──────────────────────────────────────────────────────────
\newcommand{\mkmtitlepage}{%
    \begin{titlepage}
    \begin{center}
    \vspace*{1cm}
    \begin{figure}[H]\centering
        \includegraphics[width=0.4\textwidth]{../shared/Logos/MKM_logo.png}
    \end{figure}
    \vspace{1.5cm}
    {\Huge\bfseries\textcolor{mkmblue}{\doctitle}\par}
    \vspace{0.5cm}
    {\Large\textcolor{gray}{\docsubtitle}\par}
    \vspace{2cm}
    {\large MKM Research Labs\par}
    \vspace{0.5cm}
    {\large Version \docversion{} --- \docdate\par}
    \vspace{0.5cm}
    {\large \docauthor\par}
    \vspace{2cm}
    \textcolor{red}{\textbf{CONFIDENTIAL --- PROPRIETARY}}
    \end{center}
    \end{titlepage}
}

\newcommand{\mkmlegalpage}{%
    \clearpage
    \thispagestyle{empty}
    \begin{center}
    \large\textbf{Legal Notice}
    \vspace{1em}

    \textbf{PROPRIETARY AND CONFIDENTIAL}
    \end{center}

    \noindent This document, including all algorithms, models, methodology, formulas,
    calculations, and intellectual concepts contained herein, constitutes the exclusive
    intellectual property and confidential trade secrets of MKM Research Labs (``MKM'').

    \vspace{0.5em}
    \noindent Any usage, reproduction, distribution, modification, or disclosure of this
    document or any portion thereof, without the express prior written authorisation from
    MKM Research Labs is strictly prohibited and constitutes an infringement of intellectual
    property rights.

    \vspace{0.5em}
    \noindent The document is provided on an ``as is'' basis. MKM Research Labs makes no
    representations or warranties of any kind, express or implied, regarding its accuracy,
    reliability, or suitability for any particular purpose.

    \vspace{1em}
    \begin{center}
    All rights reserved. \copyright\ 2021--2026 MKM Research Labs.

    \vspace{0.5em}
    Governed by the laws of the United Kingdom.
    \end{center}
    \clearpage
}

\newcommand{\mkmhistorypage}[1]{%
    \section{Document History}
    \begin{table}[ht]
        \centering
        \begin{tabular}{l|l|l|l}
            \textbf{Date} & \textbf{Version} & \textbf{Description} & \textbf{Author} \\
            \hline
            #1
        \end{tabular}
        \label{tab:revision_history}
    \end{table}
    \clearpage
}

% ─── Bank validation protocol section commands ──────────────────────────────
% Use these to ensure every model doc has the required sections.
% They are just styled \section calls — content is provided by each doc.

\newcommand{\mkmexecsummary}{\section{Executive Summary}}
\newcommand{\mkmpurpose}{\section{Model Purpose and Scope}}
\newcommand{\mkmframework}{\section{Mathematical Framework}}
\newcommand{\mkminputs}{\section{Input Parameters and Calibration}}
\newcommand{\mkmimplementation}{\section{Implementation Details}}
\newcommand{\mkmsensitivity}{\section{Sensitivity Analysis}}
\newcommand{\mkmvalidation}{\section{Validation and Backtesting}}
\newcommand{\mkmtestresults}{\subsection{Automated Test Results}% Auto-generated by test_results generator
% Generated: 2026-03-18 08:10:58
% Do not edit manually --- regenerate with: python -m docs.models.test_results.generator

\begin{table}[H]
    \centering
    \caption{Test Results Summary --- Flood Probability Classifier (MKM-FC-001)}
    \label{tab:test_summary_mkm_fc_001}
    \begin{tabular}{lr}
        \toprule
        Total Tests & 34 \\
        \textcolor{green!60!black}{Passed} & 34 \\
        \textcolor{red}{Failed} & 0 \\
        Skipped & 0 \\
        Duration & 0.34s \\
        \bottomrule
    \end{tabular}
\end{table}

\begin{longtable}{p{6cm}p{5.5cm}cc}
    \caption{Detailed Test Results --- Flood Probability Classifier} \\
    \label{tab:test_detail_mkm_fc_001}
    \toprule
    \textbf{Test Description} & \textbf{Acceptance Criteria} & \textbf{Result} & \textbf{Time} \\
    \midrule
    \endfirsthead
    \toprule
    \textbf{Test Description} & \textbf{Acceptance Criteria} & \textbf{Result} & \textbf{Time} \\
    \midrule
    \endhead
    \bottomrule
    \endfoot
    Each storm has peak\_level\_m. & {\scriptsize \texttt{'peak\_level\_m' in s}; \texttt{s['peak\_level\_m'] > 0}} & \textcolor{green!60!black}{Pass} & 0.001s \\
    Storms have required identification fields. & {\scriptsize \texttt{'storm\_id' in s}; \texttt{'name' in s} (+3 more)} & \textcolor{green!60!black}{Pass} & 0.002s \\
    Returns 404 when stress\_storms.json missing. & {\scriptsize \texttt{resp.status\_code == 404}} & \textcolor{green!60!black}{Pass} & 0.001s \\
    Missing gauge\_id returns total count (used by startup preloader). & {\scriptsize \texttt{resp.status\_code == 200}; \texttt{data['status'] == 'success'} (+2 more)} & \textcolor{green!60!black}{Pass} & 0.002s \\
    Returns empty list for gauge with no storms. & {\scriptsize \texttt{resp.status\_code == 200}; \texttt{data['storms'] == []}} & \textcolor{green!60!black}{Pass} & 0.002s \\
    Returns storm list for specified gauge. & {\scriptsize \texttt{resp.status\_code == 200}; \texttt{data['status'] == 'success'} (+2 more)} & \textcolor{green!60!black}{Pass} & 0.002s \\
    Storms sorted by peak level descending (worst first). & {\scriptsize \texttt{peaks == sorted(peaks, reverse=True)}} & \textcolor{green!60!black}{Pass} & 0.002s \\
    Each gauge has gauge\_id, gauge\_name, lon, lat, trade\_count. & {\scriptsize \texttt{'gauge\_id' in g}; \texttt{'gauge\_name' in g} (+3 more)} & \textcolor{green!60!black}{Pass} & 0.013s \\
    Gauges with trades show non-zero trade\_count. & {\scriptsize \texttt{sum(counts) > 0}} & \textcolor{green!60!black}{Pass} & 0.017s \\
    Returns gauge list with trade counts. & {\scriptsize \texttt{resp.status\_code == 200}; \texttt{data['status'] == 'success'} (+2 more)} & \textcolor{green!60!black}{Pass} & 0.014s \\
    Gauges sorted west to east (ascending longitude). & {\scriptsize \texttt{lons == sorted(lons)}} & \textcolor{green!60!black}{Pass} & 0.013s \\
    Each hour has water\_level, p\_flood, portfolio fields. & {\scriptsize \texttt{'hour' in h}; \texttt{'water\_level' in h} (+4 more)} & \textcolor{green!60!black}{Pass} & 0.023s \\
    Non-existent storm returns 404. & {\scriptsize \texttt{resp.status\_code == 404}} & \textcolor{green!60!black}{Pass} & 0.003s \\
    Severe storm triggers knock-out (P(flood)=100\%). & {\scriptsize \texttt{data['hourly'][first\_breach]['p\_flood'] == 1.0}; \texttt{h['p\_flood'] is None}} & \textcolor{green!60!black}{Pass} & 0.022s \\
    Missing gauge\_id or storm\_id returns 400. & {\scriptsize \texttt{resp.status\_code == 400}} & \textcolor{green!60!black}{Pass} & 0.000s \\
    No JSON body returns error. & {\scriptsize \texttt{resp.status\_code in (400, 500)}} & \textcolor{green!60!black}{Pass} & 0.002s \\
    Response includes portfolio-level stress summary. & {\scriptsize \texttt{'num\_trades' in summary}; \texttt{'total\_notional' in summary} (+5 more)} & \textcolor{green!60!black}{Pass} & 0.023s \\
    Response includes probability surface grid. & {\scriptsize \texttt{'probability\_surface' in data}; \texttt{'water\_levels' in surface} (+7 more)} & \textcolor{green!60!black}{Pass} & 0.024s \\
    POST returns 60-hour forecast array. & {\scriptsize \texttt{resp.status\_code == 200}; \texttt{data['status'] == 'success'} (+1 more)} & \textcolor{green!60!black}{Pass} & 0.021s \\
    Response includes per-trade summary. & {\scriptsize \texttt{'trades' in data}; \texttt{len(data['trades']) == 2} (+4 more)} & \textcolor{green!60!black}{Pass} & 0.021s \\
    Surface water levels cap at severe — no rows above severe. & {\scriptsize \texttt{max(surface['water\_levels']) <= severe}} & \textcolor{green!60!black}{Pass} & 0.041s \\
    Surface must cover full 168h storm horizon regardless of KO. & {\scriptsize \texttt{max(surface['hours']) >= 164}} & \textcolor{green!60!black}{Pass} & 0.033s \\
    Empty gauge\_id → 400 (line 217). & {\scriptsize \texttt{resp.status\_code == 400}} & \textcolor{green!60!black}{Pass} & 0.001s \\
    No open trades at gauge → 404 (line 250). & {\scriptsize \texttt{resp.status\_code == 404}} & \textcolor{green!60!black}{Pass} & 0.002s \\
    Null JSON body → 400 (line 205). & {\scriptsize \texttt{resp.status\_code == 400}} & \textcolor{green!60!black}{Pass} & 0.000s \\
    Predictor None (no model dir) → 404. & {\scriptsize \texttt{resp.status\_code == 404}} & \textcolor{green!60!black}{Pass} & 0.019s \\
    Gaugets file present → hydrograph built from real readings (lines 258-279). & {\scriptsize \texttt{resp.status\_code == 200}; \texttt{data['status'] == 'success'}} & \textcolor{green!60!black}{Pass} & 0.032s \\
    \_get\_predictor returns cached value without re-importing. & {\scriptsize \texttt{result is sentinel}} & \textcolor{green!60!black}{Pass} & 0.000s \\
    \_get\_predictor returns None when both stressm and legacy stress dirs are empty. & {\scriptsize \texttt{pred is None}} & \textcolor{green!60!black}{Pass} & 0.000s \\
    stress/run returns 500 on engine error (lines 302-306). & {\scriptsize \texttt{resp.status\_code == 500}} & \textcolor{green!60!black}{Pass} & 0.003s \\
    stress/gauges returns 500 when gauge loader raises (lines 122-124). & {\scriptsize \texttt{resp.status\_code == 500}} & \textcolor{green!60!black}{Pass} & 0.001s \\
    stress/storms returns 500 when storm loader raises (lines 185-187). & {\scriptsize \texttt{resp.status\_code == 500}} & \textcolor{green!60!black}{Pass} & 0.001s \\
    peak\_position=0 uses decay-only path. & {\scriptsize \texttt{len(levels) == 10}; \texttt{all(isinstance(v, float) for v in levels)}} & \textcolor{green!60!black}{Pass} & 0.000s \\
    remaining=0 uses frac=1.0 path. & {\scriptsize \texttt{len(levels) == 5}} & \textcolor{green!60!black}{Pass} & 0.000s \\
\end{longtable}


\subsection{Test Document History}
\begin{table}[ht]
    \centering
    \begin{tabular}{l|l|l|l}
        \textbf{Date} & \textbf{Version} & \textbf{Description} & \textbf{Author} \\
        \hline
        04-Mar-2026 & 1.0 & Initial test suite (15 tests) & D. Kelly \\
        05-Mar-2026 & 1.1 & 7 test(s) added (total: 22) & D. Kelly \\
        09-Mar-2026 & 1.2 & 56 test(s) added; 22 test(s) removed (total: 56) & D. Kelly \\
        09-Mar-2026 & 1.3 & 4 test(s) added; 4 test(s) removed (total: 56) & D. Kelly \\
        10-Mar-2026 & 1.4 & 12 test(s) added (total: 68) & D. Kelly \\
        10-Mar-2026 & 1.5 & 3 test(s) added (total: 71) & D. Kelly \\
        12-Mar-2026 & 1.6 & 30 test(s) added (total: 101) & D. Kelly \\
        16-Mar-2026 & 1.7 & 37 test(s) removed (total: 64) & D. Kelly \\
        17-Mar-2026 & 1.8 & 1 test(s) added; 1 test(s) removed (total: 64) & D. Kelly \\
        18-Mar-2026 & 1.9 & 30 test(s) removed (total: 34) & D. Kelly \\
    \end{tabular}
\end{table}}
\newcommand{\mkmlimitations}{\section{Model Limitations and Known Weaknesses}}

\newcommand{\mkmgovernance}{\section{Governance Validation Questions}
\noindent The following nine questions are mandated by Chapter~5 of the
\emph{Model Risk Governance for Vendors} handbook and must be explicitly
addressed for every model in the inventory. Response status is tracked in the
Model Governance Dashboard (MRC portal). \par\medskip}

\newcommand{\mkmvqitem}[3]{\subsection{Q#1 --- #2}
\noindent\textit{``#3''}\par\medskip\noindent}

\newcommand{\mkmchangelog}{\section{Change History}}
  (before \begin{document})
%
% Required commands before \input:
%   \newcommand{\doctitle}{...}
%   \newcommand{\docsubtitle}{...}
%   \newcommand{\docversion}{1.0}
%   \newcommand{\docdate}{13-February-2026}
%   \newcommand{\docauthor}{David K Kelly}

% ─── Packages ───────────────────────────────────────────────────────────────
\usepackage[utf8]{inputenc}
\usepackage{amsmath,amsfonts,amssymb}
\usepackage{listings}
\usepackage{color,xcolor}
\usepackage{graphicx}
\usepackage{hyperref}
\usepackage[toc,page]{appendix}
\usepackage{tabto}
\usepackage{float}
\usepackage{comment}
\usepackage[en-GB]{datetime2}
\usepackage{setspace}
\usepackage{subcaption}
\usepackage{booktabs}
\usepackage{longtable}
\usepackage{adjustbox}
\usepackage{enumitem}
\usepackage{fancyhdr}
\usepackage{lastpage}

% ─── Page geometry ──────────────────────────────────────────────────────────
\onehalfspacing
\usepackage{geometry}
\geometry{a4paper, total={170mm,247mm}, left=20mm, top=25mm, bottom=25mm}

% ─── Headers and footers ───────────────────────────────────────────────────
\pagestyle{fancy}
\fancyhf{}
\fancyhead[L]{\small\textcolor{gray}{MKM Research Labs}}
\fancyhead[R]{\small\textcolor{gray}{\doctitle{} --- v\docversion}}
\fancyfoot[C]{\small\textcolor{gray}{CONFIDENTIAL}}
\fancyfoot[R]{\small\textcolor{gray}{Page \thepage\ of \pageref{LastPage}}}
\renewcommand{\headrulewidth}{0.4pt}
\renewcommand{\footrulewidth}{0.4pt}

% ─── Colours ────────────────────────────────────────────────────────────────
\definecolor{mkmblue}{HTML}{1565C0}
\definecolor{codegreen}{rgb}{0,0.6,0}
\definecolor{codegray}{rgb}{0.5,0.5,0.5}
\definecolor{codepurple}{rgb}{0.58,0,0.82}
\definecolor{backcolour}{rgb}{0.95,0.95,0.92}

% ─── Hyperref setup ─────────────────────────────────────────────────────────
\hypersetup{
    colorlinks=true,
    linkcolor=mkmblue,
    citecolor=mkmblue,
    urlcolor=mkmblue
}

% ─── Code listings ──────────────────────────────────────────────────────────
\lstdefinestyle{mystyle}{
    backgroundcolor=\color{backcolour},
    commentstyle=\color{codegreen},
    keywordstyle=\color{magenta},
    numberstyle=\tiny\color{codegray},
    stringstyle=\color{codepurple},
    basicstyle=\ttfamily\footnotesize,
    breakatwhitespace=false,
    breaklines=true,
    captionpos=b,
    keepspaces=true,
    numbers=left,
    numbersep=5pt,
    showspaces=false,
    showstringspaces=false,
    showtabs=false,
    tabsize=2
}
\lstset{style=mystyle}

% ─── Section formatting ────────────────────────────────────────────────────
\setcounter{tocdepth}{3}
\setcounter{secnumdepth}{3}

% ─── Shared macros ──────────────────────────────────────────────────────────
\newcommand{\mkmtitlepage}{%
    \begin{titlepage}
    \begin{center}
    \vspace*{1cm}
    \begin{figure}[H]\centering
        \includegraphics[width=0.4\textwidth]{../shared/Logos/MKM_logo.png}
    \end{figure}
    \vspace{1.5cm}
    {\Huge\bfseries\textcolor{mkmblue}{\doctitle}\par}
    \vspace{0.5cm}
    {\Large\textcolor{gray}{\docsubtitle}\par}
    \vspace{2cm}
    {\large MKM Research Labs\par}
    \vspace{0.5cm}
    {\large Version \docversion{} --- \docdate\par}
    \vspace{0.5cm}
    {\large \docauthor\par}
    \vspace{2cm}
    \textcolor{red}{\textbf{CONFIDENTIAL --- PROPRIETARY}}
    \end{center}
    \end{titlepage}
}

\newcommand{\mkmlegalpage}{%
    \clearpage
    \thispagestyle{empty}
    \begin{center}
    \large\textbf{Legal Notice}
    \vspace{1em}

    \textbf{PROPRIETARY AND CONFIDENTIAL}
    \end{center}

    \noindent This document, including all algorithms, models, methodology, formulas,
    calculations, and intellectual concepts contained herein, constitutes the exclusive
    intellectual property and confidential trade secrets of MKM Research Labs (``MKM'').

    \vspace{0.5em}
    \noindent Any usage, reproduction, distribution, modification, or disclosure of this
    document or any portion thereof, without the express prior written authorisation from
    MKM Research Labs is strictly prohibited and constitutes an infringement of intellectual
    property rights.

    \vspace{0.5em}
    \noindent The document is provided on an ``as is'' basis. MKM Research Labs makes no
    representations or warranties of any kind, express or implied, regarding its accuracy,
    reliability, or suitability for any particular purpose.

    \vspace{1em}
    \begin{center}
    All rights reserved. \copyright\ 2021--2026 MKM Research Labs.

    \vspace{0.5em}
    Governed by the laws of the United Kingdom.
    \end{center}
    \clearpage
}

\newcommand{\mkmhistorypage}[1]{%
    \section{Document History}
    \begin{table}[ht]
        \centering
        \begin{tabular}{l|l|l|l}
            \textbf{Date} & \textbf{Version} & \textbf{Description} & \textbf{Author} \\
            \hline
            #1
        \end{tabular}
        \label{tab:revision_history}
    \end{table}
    \clearpage
}

% ─── Bank validation protocol section commands ──────────────────────────────
% Use these to ensure every model doc has the required sections.
% They are just styled \section calls — content is provided by each doc.

\newcommand{\mkmexecsummary}{\section{Executive Summary}}
\newcommand{\mkmpurpose}{\section{Model Purpose and Scope}}
\newcommand{\mkmframework}{\section{Mathematical Framework}}
\newcommand{\mkminputs}{\section{Input Parameters and Calibration}}
\newcommand{\mkmimplementation}{\section{Implementation Details}}
\newcommand{\mkmsensitivity}{\section{Sensitivity Analysis}}
\newcommand{\mkmvalidation}{\section{Validation and Backtesting}}
\newcommand{\mkmtestresults}{\subsection{Automated Test Results}% Auto-generated by test_results generator
% Generated: 2026-03-18 08:10:58
% Do not edit manually --- regenerate with: python -m docs.models.test_results.generator

\begin{table}[H]
    \centering
    \caption{Test Results Summary --- Flood Probability Classifier (MKM-FC-001)}
    \label{tab:test_summary_mkm_fc_001}
    \begin{tabular}{lr}
        \toprule
        Total Tests & 34 \\
        \textcolor{green!60!black}{Passed} & 34 \\
        \textcolor{red}{Failed} & 0 \\
        Skipped & 0 \\
        Duration & 0.34s \\
        \bottomrule
    \end{tabular}
\end{table}

\begin{longtable}{p{6cm}p{5.5cm}cc}
    \caption{Detailed Test Results --- Flood Probability Classifier} \\
    \label{tab:test_detail_mkm_fc_001}
    \toprule
    \textbf{Test Description} & \textbf{Acceptance Criteria} & \textbf{Result} & \textbf{Time} \\
    \midrule
    \endfirsthead
    \toprule
    \textbf{Test Description} & \textbf{Acceptance Criteria} & \textbf{Result} & \textbf{Time} \\
    \midrule
    \endhead
    \bottomrule
    \endfoot
    Each storm has peak\_level\_m. & {\scriptsize \texttt{'peak\_level\_m' in s}; \texttt{s['peak\_level\_m'] > 0}} & \textcolor{green!60!black}{Pass} & 0.001s \\
    Storms have required identification fields. & {\scriptsize \texttt{'storm\_id' in s}; \texttt{'name' in s} (+3 more)} & \textcolor{green!60!black}{Pass} & 0.002s \\
    Returns 404 when stress\_storms.json missing. & {\scriptsize \texttt{resp.status\_code == 404}} & \textcolor{green!60!black}{Pass} & 0.001s \\
    Missing gauge\_id returns total count (used by startup preloader). & {\scriptsize \texttt{resp.status\_code == 200}; \texttt{data['status'] == 'success'} (+2 more)} & \textcolor{green!60!black}{Pass} & 0.002s \\
    Returns empty list for gauge with no storms. & {\scriptsize \texttt{resp.status\_code == 200}; \texttt{data['storms'] == []}} & \textcolor{green!60!black}{Pass} & 0.002s \\
    Returns storm list for specified gauge. & {\scriptsize \texttt{resp.status\_code == 200}; \texttt{data['status'] == 'success'} (+2 more)} & \textcolor{green!60!black}{Pass} & 0.002s \\
    Storms sorted by peak level descending (worst first). & {\scriptsize \texttt{peaks == sorted(peaks, reverse=True)}} & \textcolor{green!60!black}{Pass} & 0.002s \\
    Each gauge has gauge\_id, gauge\_name, lon, lat, trade\_count. & {\scriptsize \texttt{'gauge\_id' in g}; \texttt{'gauge\_name' in g} (+3 more)} & \textcolor{green!60!black}{Pass} & 0.013s \\
    Gauges with trades show non-zero trade\_count. & {\scriptsize \texttt{sum(counts) > 0}} & \textcolor{green!60!black}{Pass} & 0.017s \\
    Returns gauge list with trade counts. & {\scriptsize \texttt{resp.status\_code == 200}; \texttt{data['status'] == 'success'} (+2 more)} & \textcolor{green!60!black}{Pass} & 0.014s \\
    Gauges sorted west to east (ascending longitude). & {\scriptsize \texttt{lons == sorted(lons)}} & \textcolor{green!60!black}{Pass} & 0.013s \\
    Each hour has water\_level, p\_flood, portfolio fields. & {\scriptsize \texttt{'hour' in h}; \texttt{'water\_level' in h} (+4 more)} & \textcolor{green!60!black}{Pass} & 0.023s \\
    Non-existent storm returns 404. & {\scriptsize \texttt{resp.status\_code == 404}} & \textcolor{green!60!black}{Pass} & 0.003s \\
    Severe storm triggers knock-out (P(flood)=100\%). & {\scriptsize \texttt{data['hourly'][first\_breach]['p\_flood'] == 1.0}; \texttt{h['p\_flood'] is None}} & \textcolor{green!60!black}{Pass} & 0.022s \\
    Missing gauge\_id or storm\_id returns 400. & {\scriptsize \texttt{resp.status\_code == 400}} & \textcolor{green!60!black}{Pass} & 0.000s \\
    No JSON body returns error. & {\scriptsize \texttt{resp.status\_code in (400, 500)}} & \textcolor{green!60!black}{Pass} & 0.002s \\
    Response includes portfolio-level stress summary. & {\scriptsize \texttt{'num\_trades' in summary}; \texttt{'total\_notional' in summary} (+5 more)} & \textcolor{green!60!black}{Pass} & 0.023s \\
    Response includes probability surface grid. & {\scriptsize \texttt{'probability\_surface' in data}; \texttt{'water\_levels' in surface} (+7 more)} & \textcolor{green!60!black}{Pass} & 0.024s \\
    POST returns 60-hour forecast array. & {\scriptsize \texttt{resp.status\_code == 200}; \texttt{data['status'] == 'success'} (+1 more)} & \textcolor{green!60!black}{Pass} & 0.021s \\
    Response includes per-trade summary. & {\scriptsize \texttt{'trades' in data}; \texttt{len(data['trades']) == 2} (+4 more)} & \textcolor{green!60!black}{Pass} & 0.021s \\
    Surface water levels cap at severe — no rows above severe. & {\scriptsize \texttt{max(surface['water\_levels']) <= severe}} & \textcolor{green!60!black}{Pass} & 0.041s \\
    Surface must cover full 168h storm horizon regardless of KO. & {\scriptsize \texttt{max(surface['hours']) >= 164}} & \textcolor{green!60!black}{Pass} & 0.033s \\
    Empty gauge\_id → 400 (line 217). & {\scriptsize \texttt{resp.status\_code == 400}} & \textcolor{green!60!black}{Pass} & 0.001s \\
    No open trades at gauge → 404 (line 250). & {\scriptsize \texttt{resp.status\_code == 404}} & \textcolor{green!60!black}{Pass} & 0.002s \\
    Null JSON body → 400 (line 205). & {\scriptsize \texttt{resp.status\_code == 400}} & \textcolor{green!60!black}{Pass} & 0.000s \\
    Predictor None (no model dir) → 404. & {\scriptsize \texttt{resp.status\_code == 404}} & \textcolor{green!60!black}{Pass} & 0.019s \\
    Gaugets file present → hydrograph built from real readings (lines 258-279). & {\scriptsize \texttt{resp.status\_code == 200}; \texttt{data['status'] == 'success'}} & \textcolor{green!60!black}{Pass} & 0.032s \\
    \_get\_predictor returns cached value without re-importing. & {\scriptsize \texttt{result is sentinel}} & \textcolor{green!60!black}{Pass} & 0.000s \\
    \_get\_predictor returns None when both stressm and legacy stress dirs are empty. & {\scriptsize \texttt{pred is None}} & \textcolor{green!60!black}{Pass} & 0.000s \\
    stress/run returns 500 on engine error (lines 302-306). & {\scriptsize \texttt{resp.status\_code == 500}} & \textcolor{green!60!black}{Pass} & 0.003s \\
    stress/gauges returns 500 when gauge loader raises (lines 122-124). & {\scriptsize \texttt{resp.status\_code == 500}} & \textcolor{green!60!black}{Pass} & 0.001s \\
    stress/storms returns 500 when storm loader raises (lines 185-187). & {\scriptsize \texttt{resp.status\_code == 500}} & \textcolor{green!60!black}{Pass} & 0.001s \\
    peak\_position=0 uses decay-only path. & {\scriptsize \texttt{len(levels) == 10}; \texttt{all(isinstance(v, float) for v in levels)}} & \textcolor{green!60!black}{Pass} & 0.000s \\
    remaining=0 uses frac=1.0 path. & {\scriptsize \texttt{len(levels) == 5}} & \textcolor{green!60!black}{Pass} & 0.000s \\
\end{longtable}


\subsection{Test Document History}
\begin{table}[ht]
    \centering
    \begin{tabular}{l|l|l|l}
        \textbf{Date} & \textbf{Version} & \textbf{Description} & \textbf{Author} \\
        \hline
        04-Mar-2026 & 1.0 & Initial test suite (15 tests) & D. Kelly \\
        05-Mar-2026 & 1.1 & 7 test(s) added (total: 22) & D. Kelly \\
        09-Mar-2026 & 1.2 & 56 test(s) added; 22 test(s) removed (total: 56) & D. Kelly \\
        09-Mar-2026 & 1.3 & 4 test(s) added; 4 test(s) removed (total: 56) & D. Kelly \\
        10-Mar-2026 & 1.4 & 12 test(s) added (total: 68) & D. Kelly \\
        10-Mar-2026 & 1.5 & 3 test(s) added (total: 71) & D. Kelly \\
        12-Mar-2026 & 1.6 & 30 test(s) added (total: 101) & D. Kelly \\
        16-Mar-2026 & 1.7 & 37 test(s) removed (total: 64) & D. Kelly \\
        17-Mar-2026 & 1.8 & 1 test(s) added; 1 test(s) removed (total: 64) & D. Kelly \\
        18-Mar-2026 & 1.9 & 30 test(s) removed (total: 34) & D. Kelly \\
    \end{tabular}
\end{table}}
\newcommand{\mkmlimitations}{\section{Model Limitations and Known Weaknesses}}

\newcommand{\mkmgovernance}{\section{Governance Validation Questions}
\noindent The following nine questions are mandated by Chapter~5 of the
\emph{Model Risk Governance for Vendors} handbook and must be explicitly
addressed for every model in the inventory. Response status is tracked in the
Model Governance Dashboard (MRC portal). \par\medskip}

\newcommand{\mkmvqitem}[3]{\subsection{Q#1 --- #2}
\noindent\textit{``#3''}\par\medskip\noindent}

\newcommand{\mkmchangelog}{\section{Change History}}

\begin{document}
\mkmtitlepage
\mkmlegalpage
\tableofcontents
\clearpage

% ===================================================================
% 1. EXECUTIVE SUMMARY
% ===================================================================
\mkmexecsummary

This document presents the mathematical framework, implementation, and validation of
the Flood Probability Classifier (MKM-FC-001), a Gradient Boosting Machine (GBM) model
that estimates per-gauge flood probability as a function of log-transformed water level,
storm time, and hydrograph derivatives:
$F\!\left(\ln\tfrac{w}{s},\, \ln\tfrac{t{+}1}{T},\, \Delta\ln\tfrac{w}{s},\,
\Delta^2\ln\tfrac{w}{s}\right) \to P(\text{flood})$,
where $s$ is the severe flood threshold and $T = 168$ hours.

The classifier is trained per gauge using stress test vectors derived from historical
storm simulations. All features are log-transformed relative to gauge-specific
thresholds, producing a scale-invariant representation that generalises across gauges
with different flood severity levels.

\textbf{v1.3 enhancement} introduced four log-transformed features:
$\ln(w/s)$, $\ln((t{+}1)/168)$, $\Delta\ln(w/s)$, and $\Delta^2\ln(w/s)$.
This produced a mean AUC of 0.9525 across 40 gauges, a 3.2\% improvement over v1.2.

\textbf{v2.0 upgrade (multi-storm):} Version~2.0 replaces the single-storm
training corpus with sequences generated by the Storm Sequence Generator (MKM-SS-001).
Instead of 1{,}000 independent storms ($\times$168h = 168{,}000 vectors per gauge),
the classifier now trains on \textbf{10{,}000 multi-storm sequences}
($\times$168h = \textbf{1{,}680{,}000 vectors per gauge})---ten times the data.
Crucially, multi-storm hydrographs contain \emph{multi-peak} shapes driven by
compounding antecedent-wetness effects that single-storm data cannot represent.
The feature space and hyperparameters are unchanged from v1.3; the gains come
entirely from richer, larger training data.

The mean AUC-ROC across all 52 Thames gauges improved from 0.9525 (v1.3) to
\textbf{$\approx$0.993} (v2.0), a further 4.2\% improvement.
The stressm classifiers are stored in \texttt{data/output/stressm/}
and trained via \texttt{phys.py port --stressm --train-classifier}.

\textbf{Key risk:} The model is trained on synthetically noised data (time-decaying
Gaussian perturbation), which creates intentional early-hour uncertainty. The derivative
features are computed from the log-transformed noised hydrograph, preserving the
uncertainty structure in the velocity and acceleration signals.

% ===================================================================
% 2. MODEL PURPOSE AND SCOPE
% ===================================================================
\mkmpurpose

\subsection{Purpose}
The Flood Probability Classifier serves two functions within the platform:

\begin{itemize}[nosep]
    \item \textbf{Stress testing:} Given a stress storm scenario, estimate the
          probability of flood alert breach at each gauge as a function of
          evolving water levels during the storm event.
    \item \textbf{P\&L estimation:} Convert flood probabilities into cash-equivalent
          stress P\&L for PRS trade positions, using the CDS-in-stress analogy:
          $\text{cash\_price} = \text{notional} \times P(\text{flood}) \times \text{direction}$.
\end{itemize}

\subsection{Scope}
\begin{itemize}[nosep]
    \item \textbf{In scope:} Per-gauge binary classification (flood/no-flood) using
          water level, temporal, and hydrograph derivative features.
    \item \textbf{Out of scope:} Multi-gauge spatial correlation, flood severity
          estimation (depth/velocity), and climate change projection.
\end{itemize}

\subsection{Model Chain Position}
\label{sec:model_chain}
The classifier now sits downstream of the Storm Sequence Generator:
\[
\text{MKM-SI-001} \to
\textbf{MKM-SS-001}\ (\text{sequences})
\to \text{MKM-SG-001}\ (\text{hydrograph})
\to \textbf{MKM-FC-001}\ (\text{v2.0})
\to \text{Stress P\&L}
\]
For single-storm stress testing (Trading Desk Tab~7), the v1.3 classifiers stored in
\texttt{data/output/stress/} remain in use. The v2.0 stressm classifiers stored in
\texttt{data/output/stressm/} power the Portfolio Stress tab (Tab~8).

% ===================================================================
% 3. MATHEMATICAL FRAMEWORK
% ===================================================================
\mkmframework

\subsection{Model Family}
The model uses \texttt{sklearn.ensemble.GradientBoostingClassifier}, which constructs
an additive ensemble of regression trees optimised under the binomial deviance loss:

\[
\mathcal{L}(y, F) = -\sum_{i=1}^{N} \left[ y_i \log p_i + (1 - y_i) \log(1 - p_i) \right]
\]

where $p_i = \sigma(F(x_i))$ and $\sigma$ is the logistic function. The model is built
sequentially:

\[
F_m(x) = F_{m-1}(x) + \nu \cdot h_m(x)
\]

where $h_m$ is a regression tree fit to the negative gradient of the loss, and $\nu$
is the learning rate.

\subsection{Feature Space}
The model uses a four-dimensional log-transformed feature space:

\begin{center}
\begin{tabular}{llp{7.5cm}}
\toprule
\textbf{Feature} & \textbf{Type} & \textbf{Description} \\
\midrule
\texttt{log\_h\_s} & Continuous &
    $\ln\!\left(\tilde{w}(t) / s\right)$; water level relative to severe threshold $s$.
    Zero at severe level, negative below, positive above. \\
\texttt{log\_t\_end} & Continuous &
    $\ln\!\left((t{+}1) / 168\right)$; log storm time.
    Ranges from $\approx{-5.1}$ (hour 0) to $0$ (hour 167). \\
\texttt{delta\_log\_h} & Continuous &
    $\Delta\ln(w/s)(t) = \texttt{log\_h\_s}(t) - \texttt{log\_h\_s}(t{-}1)$;
    log velocity (positive rising, negative falling). \\
\texttt{delta2\_log\_h} & Continuous &
    $\Delta^2\ln(w/s)(t) = \texttt{delta\_log\_h}(t) - \texttt{delta\_log\_h}(t{-}1)$;
    log acceleration. \\
\bottomrule
\end{tabular}
\end{center}

\noindent All derivative features are computed from the log-transformed noised
hydrograph $\ln(\tilde{w}(t)/s)$, not the clean signal. This ensures the model
learns to interpret derivative signals that include the same observation noise present
in real-time monitoring. At prediction time, raw derivatives $\Delta w$ and
$\Delta^2 w$ are passed in and converted to log-space internally by the predictor.

\subsection{Physical Motivation for Log Transformation}
The raw-level feature space has two limitations:

\begin{enumerate}[nosep]
    \item \textbf{Gauge heterogeneity:} Thames gauges have severe thresholds ranging
          from approximately 3\,m to 7\,m AOD. A water level of 4.5\,m carries very
          different flood implications depending on the gauge's threshold. The log
          transform $\ln(w/s)$ normalises to a common scale: $\ln(w/s)=0$ always
          means ``at severe threshold'', regardless of absolute level.
    \item \textbf{Velocity interpretation:} A rise of 0.3\,m/h from 1\,m is alarming;
          the same rise from 5\,m is less informative. Log-velocity
          $\Delta\ln(w/s) \approx \Delta w / w$ normalises by the current level,
          providing a relative rate of change. This is analogous to using log-returns
          in financial risk models.
\end{enumerate}

\noindent The time feature $\ln((t{+}1)/168)$ is restored from v1.2 in log scale.
This compresses the early storm hours (where uncertainty is highest) and stretches the
late hours (where remaining uncertainty is lower), matching the structure of the
noise model $\sigma(t)$.

\subsection{Rising vs Falling Water}
The log-velocity feature $\Delta\ln(w/s)$ preserves the directional asymmetry
introduced in v1.2:

\begin{itemize}[nosep]
    \item \textbf{Rising water} ($\Delta\ln(w/s) > 0$): The storm has not yet peaked.
          $P(\text{flood})$ should reflect uncertainty about the eventual peak.
    \item \textbf{Falling water} ($\Delta\ln(w/s) < 0$): The storm has already peaked.
          If near alert level, the peak was likely above it.
          $P(\text{flood})$ should be higher, reflecting stronger evidence.
\end{itemize}

\noindent The log-acceleration $\Delta^2\ln(w/s)$ further discriminates between a
decelerating rise (approaching peak) and an accelerating rise (early surge).

\subsection{Target Variable}
Binary label $y \in \{0, 1\}$ indicating whether the storm breaches the gauge's
flood alert trigger level at any point during the event. The label is computed from
the stress test vector:

\[
y_i = \begin{cases}
1 & \text{if } \max_{t} w_i(t) \geq \text{alert\_level} \\
0 & \text{otherwise}
\end{cases}
\]

\subsection{Training Data}

\subsubsection{v2.0 (stressm): Multi-Storm Sequence Training}

Version~2.0 trains on the full catalogue of 10{,}000 \texttt{StormSequence} objects
generated by MKM-SS-001 (Section~\ref{sec:model_chain}). Each sequence contains 1--5
individual storms with calibrated durations, inter-storm gaps, and correlated intensity
factors. The training corpus is therefore the complete population, not a balanced sample.

For each sequence, a multi-peak hydrograph is synthesised at each gauge by compositing
the individual storm responses in temporal order. Each storm contributes a water level
arc derived from its intensity factor and the gauge's spatial response parameters
(\texttt{base\_level\_m}, \texttt{level\_change\_m}, \texttt{peak\_position}).
Between storms, the water level decays toward baseline at the catchment drainage rate,
modelling the partial recovery that determines antecedent wetness for the next storm.

The same time-decaying Gaussian noise is applied per storm arc:
\[
\tilde{w}(t) = w(t) + \varepsilon(t), \quad
\varepsilon(t) \sim \mathcal{N}\!\left(0, \sigma(t)^2\right)
\]
\[
\sigma(t) = 0.5 \cdot e^{-3t / 167} + 0.1
\]

The feature extraction is identical to v1.3; the features are computed from the full
168-hour noised multi-peak hydrograph:
\[
\mathbf{x}_i = \left[\, \ell(t),\; \tau(t),\; \Delta\ell(t),\; \Delta^2\ell(t) \,\right]
\in \mathbb{R}^4
\]

\noindent Total training corpus: \textbf{1{,}680{,}000 vectors per gauge}
(10{,}000 sequences $\times$ 168 hours).

\subsubsection{v1.3 (stress): Single-Storm Balanced Training}

For reference, the v1.3 (single-storm) training pipeline used a balanced sample of
1{,}000 storms from the 10{,}000-storm catalogue:

\begin{itemize}[nosep]
    \item \textbf{Alert storms (260):} All storms that breached the flood alert trigger,
          sourced from \texttt{stress\_storms.json}.
    \item \textbf{Non-alert storms (740):} Randomly sampled from the non-alert population.
\end{itemize}

Total: 168{,}000 vectors per gauge (1{,}000 storms $\times$ 168 hours).

\textbf{Why v2.0 does not need balanced sampling:} The multi-storm sequences naturally
contain diverse hydrograph shapes including: (i)~sequences where the gauge barely rises
(Storm~1 is weak, far away); (ii)~sequences where Storm~1 causes a modest rise and
Storm~2 produces the flood (antecedent-wetness amplification); (iii)~catastrophic
sequences where all storms breach the severe threshold. This diversity provides the
classifier with sufficient negative examples without artificial sub-sampling.

Log-level features are computed identically to v1.3:
\[
\ell(t) = \ln\!\left(\max\!\left(\frac{\tilde{w}(t)}{s},\; 10^{-8}\right)\right),
\quad
\Delta\ell(t) = \ell(t) - \ell(t{-}1), \quad
\Delta^2\ell(t) = \Delta\ell(t) - \Delta\ell(t{-}1)
\]
\[
\tau(t) = \ln\!\left(\frac{t + 1}{168}\right)
\]
with boundary conditions $\Delta\ell(0) = \Delta^2\ell(0) = 0$.

% ===================================================================
% 4. HYPERPARAMETERS
% ===================================================================
\section{Hyperparameters}
\label{sec:hyperparameters}

\begin{center}
\begin{tabular}{lrl}
\toprule
\textbf{Parameter} & \textbf{Value} & \textbf{Rationale} \\
\midrule
\texttt{n\_estimators} & 300 & Sufficient for convergence without overfitting \\
\texttt{max\_depth} & 3 & Captures water-level $\times$ hour interaction \\
\texttt{learning\_rate} & 0.05 & Shrinkage paired with 300 rounds \\
\texttt{subsample} & 0.8 & Row sampling to reduce variance \\
\texttt{min\_samples\_leaf} & 50 & Prevents small-node overfitting \\
\texttt{random\_state} & 42 & Reproducibility \\
\texttt{test\_size} & 0.2 & 80/20 train-test split \\
\bottomrule
\end{tabular}
\end{center}

The stressm trainer holds out 20\% of storms, stratified on whether the storm floods,
so every storm is entirely in training or entirely in test
(Section~\ref{sec:leakage}).

% ===================================================================
% 5. IMPLEMENTATION
% ===================================================================
\section{Implementation Details}

\subsection{Source Code}
\begin{itemize}[nosep]
    \item \textbf{Primary module:} \texttt{src/models/stress/flood\_classifier.py}
    \item \textbf{Training class:} \texttt{FloodClassifierTrainer}
    \item \textbf{Prediction class:} \texttt{FloodPredictor}
\end{itemize}

\subsection{Training Pipelines}

\subsubsection{v2.0 stressm Pipeline}
\begin{enumerate}[nosep]
    \item Generate 10{,}000 StormSequence objects via MKM-SS-001 (\texttt{--stressm})
    \item For each gauge, synthesise 168-hour multi-peak hydrograph per sequence
    \item Apply time-decaying noise; compute 4 log-transformed feature vectors
    \item Build feature matrix $\mathbf{X} \in \mathbb{R}^{1{,}680{,}000 \times 4}$
          and label vector $\mathbf{y} \in \{0,1\}^{1{,}680{,}000}$
    \item Skip single-class gauges
    \item 80/20 train-test split by storm with \texttt{random\_state=42};
          near-miss hydrographs stay in training
    \item Fit \texttt{GradientBoostingClassifier} on training set
    \item Evaluate on test set: accuracy, AUC-ROC, Brier score, log loss
    \item Serialise model to \texttt{data/output/stressm/GAUGE-\{id\}.joblib}
    \item Write summary to \texttt{data/output/stressm/training\_summary.json}
\end{enumerate}

\subsubsection{v1.3 stress Pipeline (retained for single-storm Tab~7)}
\begin{enumerate}[nosep]
    \item Load stress test vectors from \texttt{data/input/thames/stress\_vectors.json}
    \item For each gauge, extract feature matrix $\mathbf{X} \in \mathbb{R}^{168{,}000 \times 4}$
          and label vector $\mathbf{y} \in \{0,1\}^{168{,}000}$
    \item Stratified 80/20 split; fit; evaluate; serialise to
          \texttt{data/output/stress/GAUGE-\{id\}.joblib}
\end{enumerate}

\subsection{Prediction Interface}
Three prediction modes are provided:

\begin{itemize}[nosep]
    \item \texttt{predict(gauge\_id, water\_level, hour, delta\_w, delta2\_w)}
          --- Single-point $\hat{p}$; \texttt{delta\_w} and \texttt{delta2\_w}
          default to 0.0 (steady state) for backward compatibility
    \item \texttt{predict\_series(gauge\_id, levels, hours, delta\_ws, delta2\_ws)}
          --- Vectorised time series
    \item \texttt{predict\_surface(gauge\_id, ..., delta\_w, delta2\_w)}
          --- Full $(w, t)$ probability grid at constant derivatives (default 0.0)
\end{itemize}

\subsection{Probability Surface}
The \texttt{predict\_surface} mode produces a two-dimensional grid of flood
probabilities $P(w, t)$ over a range of water levels $w$ and storm hours $t$.
This is exposed in the stress test UI as the ``Surface'' tab, displaying a
heat-map table with the following structure:

\begin{center}
\begin{tabular}{ll}
\toprule
\textbf{Axis} & \textbf{Description} \\
\midrule
Rows (Y) & Water levels in 0.5\,m steps, descending from highest to lowest \\
Columns (X) & Storm hours at 4-hour intervals: H0, H4, H8, \ldots, H56 \\
Cells & $P(\text{flood})$ as a percentage to one decimal place \\
\bottomrule
\end{tabular}
\end{center}

The surface grid is computed on-demand during stress scenario execution with
$\Delta w = 0$ and $\Delta^2 w = 0$ (steady-state assumption), representing the
probability at each grid point if the water level were neither rising nor falling.
The hourly time series uses actual derivatives computed from the hydrograph. Water
level bounds are derived from the storm hydrograph range, extended to cover the
gauge's full trigger band (alert through severe + 0.5\,m buffer).

\textbf{Visual encoding:}
\begin{itemize}[nosep]
    \item Cells below the alert level: neutral (white) background with grey
          intensity shading proportional to $P(\text{flood})$.
    \item Cells in the alert band ($\text{alert} \leq w < \text{warning}$):
          amber background (\texttt{\#FFF8E1}).
    \item Cells in the warning band ($\text{warning} \leq w < \text{severe}$):
          orange background (\texttt{\#FFF3E0}).
    \item Cells at or above the severe level: red background (\texttt{\#FFEBEE}).
    \item Post knock-out (after the first trigger breach hour): cells are left
          blank with a grey background, reflecting contract termination.
\end{itemize}

Within each trigger band, cell opacity increases proportionally to the probability
value, creating a heat-map effect that highlights regions of high flood risk.

\subsection{Serialisation}
Models are persisted using \texttt{joblib.dump()} in the standard sklearn format.
Each gauge model is stored independently at \texttt{GAUGE-\{id\}.joblib}, enabling
selective retraining without full pipeline re-execution.

% ===================================================================
% 6. VALIDATION AND METRICS
% ===================================================================
\section{Validation and Metrics}

\subsection{Evaluation Metrics}
Four metrics are computed per gauge on the held-out test set:

\begin{center}
\begin{tabular}{lp{9cm}}
\toprule
\textbf{Metric} & \textbf{Description} \\
\midrule
Accuracy & Fraction of correct binary predictions \\
AUC-ROC & Area under the Receiver Operating Characteristic curve; measures
           discrimination ability independent of threshold \\
Brier Score & Mean squared error of probabilistic predictions; lower is better \\
Log Loss & Negative log-likelihood of predicted probabilities; penalises
           confident wrong predictions \\
\bottomrule
\end{tabular}
\end{center}

\subsection{Production Results (v2.0 stressm)}
Training across all 52 Thames catchment gauges with 10{,}000 multi-storm sequences:

\begin{itemize}[nosep]
    \item \textbf{Average AUC-ROC:} $\approx$0.993 (row-level split, before v2.2;
          see Section~\ref{sec:leakage})
    \item \textbf{Representative gauge (GAUGE-8f1d368d):} AUC 0.9927,
          accuracy 0.9776, flood rate 8.4\%
    \item \textbf{Training sequences:} 10{,}000 StormSequence objects (all types)
    \item \textbf{Vectors per gauge:} 1{,}680{,}000 (10{,}000 sequences $\times$ 168 hours)
    \item \textbf{Gauges:} 52 (full Thames portfolio; no gauges skipped)
    \item \textbf{Feature importance:} \texttt{log\_h\_s} $\approx$70\%,
          \texttt{log\_t\_end} $\approx$24\%, derivatives $\approx$6\%
\end{itemize}

\subsection{Version Comparison: v1.1 $\to$ v1.2 $\to$ v1.3 $\to$ v2.0}

The following table shows the progression across all model versions:

\begin{center}
\begin{tabular}{lccc}
\toprule
\textbf{Version} & \textbf{Key Change} & \textbf{Mean AUC} & \textbf{Vectors/gauge} \\
\midrule
v1.1 & 2 features (level, hour) & 0.9189 & 168{,}000 \\
v1.2 & + $\Delta w$, $\Delta^2 w$ (log-raw) & 0.9557 & 168{,}000 \\
v1.3 & Log-transform all features & 0.9525 & 168{,}000 \\
\textbf{v2.0} & \textbf{Multi-storm sequences; 10$\times$ data} & $\approx$\textbf{0.993} & \textbf{1{,}680{,}000} \\
\bottomrule
\end{tabular}
\end{center}

\noindent\textbf{Analysis:} The v1.1$\to$v1.2 jump (+4.0\%) came from adding derivative
features. The v1.2$\to$v1.3 slight reduction ($-$0.3\%) reflects the more rigorous
log-space normalisation (marginally harder problem, better-calibrated probabilities).
The v1.3$\to$v2.0 step (+4.2\%) comes entirely from richer, larger training data:
multi-peak hydrographs force the classifier to learn antecedent-wetness dynamics,
and 10$\times$ more data reduces tree estimation variance throughout the GBM ensemble.

\subsubsection{Asymmetry Evidence: Rising vs Falling Water}
The key diagnostic is the model's ability to distinguish rising from falling water
at the same level. The following table shows $P(\text{flood})$ for gauge
\texttt{GAUGE-32e88723} (AUC\,=\,0.957, flood rate\,=\,10.0\%) at hour\,=\,20
with $|\Delta w| = 0.4$\,m/h:

\begin{center}
\begin{tabular}{rcccc}
\toprule
\textbf{Level (m)} & \textbf{Rising} & \textbf{Steady} & \textbf{Falling}
    & \textbf{$\Delta P$} \\
\midrule
3.3 (sub-alert) & 0.003 & 0.027 & 0.103 & $-$0.099 \\
3.8 & 0.033 & 0.140 & 0.462 & $-$0.429 \\
4.3 (near alert) & 0.261 & 0.646 & 0.882 & $-$0.622 \\
4.8 (above alert) & 0.831 & 0.942 & 0.976 & $-$0.144 \\
5.3 & 0.984 & 0.992 & 0.996 & $-$0.012 \\
\bottomrule
\end{tabular}
\end{center}

\noindent The model has learned the physically correct signal:
\begin{itemize}[nosep]
    \item At sub-alert levels, \textbf{falling water implies higher $P(\text{flood})$}:
          if the water is at 4.3\,m and falling, the storm already peaked higher,
          providing strong evidence that alert was breached.
    \item At sub-alert levels, \textbf{rising water implies lower $P(\text{flood})$}:
          the storm has not yet peaked, so there is genuine uncertainty about
          whether the peak will reach the alert level.
    \item Above the alert level, all states converge to $P \approx 1$ regardless
          of direction, as the evidence for flood breach is overwhelming.
\end{itemize}

The maximum asymmetry occurs near the alert threshold (4.3\,m), where the directional
signal is most informative: $\Delta P = -0.622$. The v1.1 model would output the
``Steady'' column for both rising and falling water.

\subsection{Feature Importance}
Averaged across all 40 gauges, the four features contribute:

\begin{center}
\begin{tabular}{lr}
\toprule
\textbf{Feature} & \textbf{Importance} \\
\midrule
\texttt{water\_level} & 0.6054 \\
\texttt{hour} & 0.2233 \\
\texttt{delta\_w} & 0.1407 \\
\texttt{delta2\_w} & 0.0307 \\
\bottomrule
\end{tabular}
\end{center}

\noindent Water level remains the dominant predictor (60.5\%), as expected.
The velocity feature (\texttt{delta\_w}) captures 14.1\% of model importance---a
substantial signal that enables the rising/falling asymmetry. The acceleration
feature (\texttt{delta2\_w}) contributes 3.1\%, providing marginal refinement
for distinguishing decelerating rises (approaching peak) from accelerating surges.

\noindent Compared to v1.1, the hour feature's relative importance decreased from
$\sim$14\% to 22.3\%---however, this is a redistribution effect: in absolute terms,
the temporal signal is captured more precisely because \texttt{delta\_w} absorbs
some of the temporal variation previously attributed to \texttt{hour} alone.

\subsection{Automated Test Results}
Current unit test evidence is filed with the model record in ModelRisk (Test Results and Unit Test).

% ===================================================================
% 7. SENSITIVITY ANALYSIS
% ===================================================================
\mkmsensitivity

Sensitivity of the classifier's AUC to key hyperparameters was analysed by varying
each parameter independently around a baseline of 200 trees, depth 4 and learning
rate 0.1.  This is not the production configuration in Section~\ref{sec:hyperparameters}:
the analysis has not been rerun at production settings, because its generator has
been retired, so the findings and tables below are historical.

Key findings:
\begin{itemize}[nosep]
    \item \texttt{n\_estimators}: AUC stabilises above 100 rounds; diminishing
          returns beyond 200 (at depth 4 and learning rate 0.1).
    \item \texttt{max\_depth}: Optimal at 4; deeper trees overfit the noise.
    \item \texttt{learning\_rate}: Sweet spot at 0.1; aggressive rates ($>$0.2)
          increase variance.
\end{itemize}

% Auto-generated by sensitivity generator
% Do not edit manually --- regenerate with: python -m docs.models.sensitivities.generate_all

\begin{table}[H]
    \centering
    \caption{Classifier sensitivity to n\_estimators (gauge GAUGE-082c, max\_depth=4, lr=0.1).}
    \label{tab:sens_n_estimators}
    \begin{tabular}{rccl}
        \toprule
        \textbf{n\_estimators} & \textbf{AUC-ROC} & \textbf{Accuracy} & \textbf{} \\
        \midrule
        50 & 0.8010 & 0.7936 &  \\
        100 & 0.8002 & 0.7917 &  \\
        200 & 0.7965 & 0.7865 & \textbf{Baseline} \\
        500 & 0.7931 & 0.7824 &  \\
        \bottomrule
    \end{tabular}
\end{table}

\begin{table}[H]
    \centering
    \caption{Classifier sensitivity to max\_depth (gauge GAUGE-082c, n\_estimators=200, lr=0.1).}
    \label{tab:sens_max_depth}
    \begin{tabular}{rccl}
        \toprule
        \textbf{max\_depth} & \textbf{AUC-ROC} & \textbf{Accuracy} & \textbf{} \\
        \midrule
        2 & 0.7974 & 0.7904 &  \\
        3 & 0.7975 & 0.7878 &  \\
        4 & 0.7965 & 0.7865 & \textbf{Baseline} \\
        5 & 0.7951 & 0.7840 &  \\
        6 & 0.7928 & 0.7801 &  \\
        \bottomrule
    \end{tabular}
\end{table}

\begin{table}[H]
    \centering
    \caption{Classifier sensitivity to learning\_rate (gauge GAUGE-082c, n\_estimators=200, max\_depth=4).}
    \label{tab:sens_learning_rate}
    \begin{tabular}{rccl}
        \toprule
        \textbf{learning\_rate} & \textbf{AUC-ROC} & \textbf{Accuracy} & \textbf{} \\
        \midrule
        0.0100 & 0.7950 & 0.7869 &  \\
        0.0500 & 0.8005 & 0.7901 &  \\
        0.1000 & 0.7965 & 0.7865 & \textbf{Baseline} \\
        0.2000 & 0.7902 & 0.7811 &  \\
        0.5000 & 0.7794 & 0.7763 &  \\
        \bottomrule
    \end{tabular}
\end{table}

\begin{table}[H]
    \centering
    \caption{Classifier sensitivity to subsample fraction (gauge GAUGE-082c, n\_estimators=200, max\_depth=4, lr=0.1).}
    \label{tab:sens_subsample}
    \begin{tabular}{rccl}
        \toprule
        \textbf{subsample} & \textbf{AUC-ROC} & \textbf{Accuracy} & \textbf{} \\
        \midrule
        0.5000 & 0.7972 & 0.7891 &  \\
        0.6000 & 0.7976 & 0.7853 &  \\
        0.7000 & 0.7971 & 0.7897 &  \\
        0.8000 & 0.7965 & 0.7865 & \textbf{Baseline} \\
        0.9000 & 0.7985 & 0.7881 &  \\
        1.00 & 0.7975 & 0.7862 &  \\
        \bottomrule
    \end{tabular}
\end{table}

\begin{table}[H]
    \centering
    \caption{Feature importance at production hyperparameters --- v1.1 vs v1.2 (gauge GAUGE-082c).}
    \label{tab:sens_feature_importance}
    \begin{tabular}{lrr}
        \toprule
        \textbf{Feature} & \textbf{v1.1} & \textbf{v1.2} \\
        \midrule
        water\_level & 0.8607 & 0.6260 \\
        hour & 0.1393 & 0.2590 \\
        delta\_w & --- & 0.0960 \\
        delta2\_w & --- & 0.0190 \\
        \midrule
        AUC-ROC & 0.7965 & 0.9720 \\
        \bottomrule
    \end{tabular}
\end{table}

% =========================================================================
\section{Model Selection: Challenger Comparison}
\label{sec:challenger}
% =========================================================================

\subsection{Open Question}

\begin{quote}
\emph{Should an alternative classifier, such as a Random Forest, be
considered in place of the GBM?  What evidence shows that the GBM is
the most accurate and most stable choice for $P(\text{flood})$?}
\end{quote}

\noindent The GBM was selected on the basis of in-sample AUC.  That is
not sufficient evidence for model selection: no challenger has been
fitted on the same data, and the AUC was measured on a split that
overstated out-of-sample performance (Section~\ref{sec:leakage}).  This
section sets out the comparison that would settle the question.  The
protocol is written so that the GBM \emph{can fail}: if a challenger
wins under the pre-agreed decision rule, the challenger is adopted.

\subsection{Why Random Forest Is the Natural Challenger}

Both models are tree ensembles on the same four features, so they
differ only in how the ensemble is built:

\begin{center}
\begin{tabular}{p{3.2cm}p{5.2cm}p{5.2cm}}
\toprule
& \textbf{GBM (champion)} & \textbf{Random Forest (challenger)} \\
\midrule
Ensemble     & Sequential, each tree fits the residual of the last
             & Independent deep trees on bootstrap samples, averaged \\
Error reduced & Mainly bias & Mainly variance \\
Tuning       & Sensitive (learning rate, depth, trees)
             & Robust to defaults \\
Probabilities & Optimises log loss directly; usually well calibrated
             & Averaged votes; pulled away from 0 and 1, often needs
               recalibration \\
Seed sensitivity & Low--moderate (\texttt{subsample=0.8})
             & Low with many trees \\
\bottomrule
\end{tabular}
\end{center}

\noindent Because the classifier output is used directly as a price
($C_h = N \times p_h \times d$), calibration matters more than
ranking.  This is the main reason to expect the GBM to win, but it has
to be shown, not assumed.

\subsection{Candidate Set}

\begin{enumerate}
  \item \textbf{GBM} --- production configuration
        (\texttt{n\_estimators=300}, \texttt{max\_depth=3},
        \texttt{learning\_rate=0.05}, \texttt{subsample=0.8},
        \texttt{min\_samples\_leaf=50}).
  \item \textbf{GBM, tuned} --- the same model given the same tuning
        budget as the challengers.  If it beats the production GBM,
        the production settings, not the model family, need changing.
  \item \textbf{Random Forest} --- \texttt{n\_estimators=500},
        \texttt{max\_features=2}, \texttt{min\_samples\_leaf} tuned.
  \item \textbf{Random Forest, calibrated} --- the same, wrapped in
        isotonic or Platt calibration fitted on an inner fold.  This
        stops RF losing only because of calibration that a
        one-line fix would remove.
  \item \textbf{Logistic regression} on the four features --- a
        transparent baseline.
  \item \textbf{Flood polynomial} --- the production fallback, as the
        floor every model must beat.
\end{enumerate}

\noindent Every tree model gets the same tuning budget (for example
30 random-search trials) inside the training folds, so the GBM does
not win only because it has been tuned more.

\subsection{Prerequisite: Remove Storm Leakage}
\label{sec:leakage}

Each storm contributes 168 hourly rows, and the label
(\emph{future maximum level} $\geq w_{\text{severe}}$) is shared by
adjacent hours.  Up to version 2.1 of this document the trainer split
by row, using a stratified \texttt{train\_test\_split}, so hours from
the same storm
appeared in both training and test sets, and near-miss rows were
scored too.  The reported AUC therefore measured interpolation within
storms already seen, not prediction on new storms.  This leakage
favours deep, low-bias learners and so would flatter the Random
Forest more than the GBM.

\noindent The production trainer now holds out 20\% of storms,
stratified on whether the storm floods, keeps near-miss rows in
training and records \texttt{split} in its result.  If no such split
is possible (for example, only one storm floods) it falls back to the
row-level split and records \texttt{split} $=$ \texttt{row}.
Classifiers and metrics trained before this change keep the old
split until they are retrained, when the fitted models also change.
On one test-portfolio gauge the new split moved Brier from 0.0213 to
0.0258 and AUC from 0.988 to 0.990.  AUC can hold up because the easy
gains from leakage are offset by no longer scoring near-miss rows, so
Brier and log loss are the better guide to the change.

\noindent For the comparison itself:
\begin{itemize}
  \item Carry a \texttt{storm\_id} with every training vector.
  \item Split by storm using \texttt{GroupKFold} (5 folds), so every
        storm is entirely in train or entirely in test.
  \item Add near-miss synthetic vectors to the \emph{training} folds
        only; score on real storm vectors only.
\end{itemize}

\subsection{Accuracy Metrics}

Scored on held-out storms, per gauge and pooled:
\begin{description}[leftmargin=2em]
  \item[Primary] Brier score and log loss --- proper scoring rules
        that reward correct probabilities, not just correct ranking.
  \item[Calibration] Reliability diagram, expected calibration error,
        calibration slope and intercept.
  \item[Discrimination] AUC (secondary, kept for continuity with
        the existing thresholds).
  \item[Business] Stress P\&L error on held-out storms: the hourly
        $\sum_h N p_h d$ path compared against the realised
        knock-out outcome of that storm.
\end{description}

\subsection{Stability Tests}

``Most stable'' is tested in five ways.  Each produces a number per
model and per gauge.

\begin{enumerate}
  \item \textbf{Seed stability.}  Retrain each model with 20 random
        seeds on the same data.  Measure the standard deviation of
        $p$ on a fixed grid (the probability surface of
        Section~4) and of peak stress P\&L for a reference storm set.
  \item \textbf{Data stability.}  Retrain on 50 bootstrap resamples
        \emph{of storms} (not of rows).  Report the spread of $p$ on
        the fixed grid and the spread of the held-out Brier score.
  \item \textbf{Catalogue stability.}  Train on one generation of the
        storm catalogue and score on a regenerated catalogue
        (a new seed or new storm-control parameters).  This is the
        retraining trigger in practice, so it is the most
        decision-relevant stability test.  The benchmark tool does not
        regenerate the catalogue itself: run it once on each catalogue
        generation and compare the two result files.
  \item \textbf{Shape stability.}  $P(\text{flood})$ should not fall
        as the water level rises.  Count monotonicity violations in
        \texttt{log\_h\_s} across the grid, and measure how rough the
        finite-difference sensitivity $\partial p / \partial w$ is.
        Step-like surfaces give noisy stress Greeks.
  \item \textbf{Cross-gauge consistency.}  The distribution of
        per-gauge metrics: a model that is excellent on most gauges
        but fails badly on a few is less stable than one that is
        uniformly good.
\end{enumerate}

\subsection{Statistical Comparison}

\begin{itemize}
  \item Per gauge: paired difference in Brier score (GBM minus
        challenger) on the same held-out storms, with a 95\%
        confidence interval from a storm-level bootstrap.
  \item Across gauges: Wilcoxon signed-rank test on the per-gauge
        Brier differences.
  \item Stability metrics: report ratios (challenger dispersion /
        GBM dispersion) with bootstrap intervals.
\end{itemize}

\subsection{Decision Rule}

Agreed before the results are seen.  The GBM is retained if all of
the following hold:
\begin{enumerate}
  \item Its pooled Brier score is no worse than the best challenger's
        by more than a set tolerance (for example 2\%), and the
        Wilcoxon test does not show the challenger significantly
        better.
  \item Its seed, bootstrap and catalogue dispersion are no greater
        than the challenger's (within 0.002 in probability s.d.).
  \item It has no more monotonicity violations than the challenger
        (within 0.5 percentage points of the violation rate).
\end{enumerate}

\noindent If a challenger is better on accuracy but worse on
stability (or the reverse), the trade-off is recorded and the choice
is escalated for model-risk sign-off rather than decided by the
model owner.  If the best challenger wins on all three, it replaces
the GBM.  With fewer than six gauges the Wilcoxon test has no power, so
``significantly better'' then means every gauge's storm-bootstrap
interval favours the challenger.

\subsection{Running the Benchmark}

The comparison is implemented in the benchmark tool:
\begin{description}[leftmargin=2em]
  \item[Code] \texttt{tools/flood\_classifier\_challenger/}
  \item[Settings] \texttt{config/models/\_flood\_challenger.py}
\end{description}
\noindent It builds each gauge's training set with the production
vector builder, so it scores exactly the features and labels
production trains on:

\begin{verbatim}
python -m tools.flood_classifier_challenger \
    --gauges all
python -m tools.flood_classifier_challenger \
    --gauges G1,G2 --seeds 20 --bootstraps 50
python -m tools.flood_classifier_challenger \
    --synthetic 3     # no storm data needed
\end{verbatim}

\noindent Defaults are 1{,}000 storms per gauge, 5 storm folds,
6 tuning trials, 10 seeds and 20 storm bootstraps; a sign-off run
uses 20 seeds and 50 bootstraps as above.  Fits run in parallel on
every core.  The run writes \texttt{challenger/challenger\_results.json}
(every metric, the tuned parameters per fold and the probability
surfaces) and the table below.

\subsection{Results}

\IfFileExists{challenger/challenger_results.tex}{%
  % Generated by tools/flood_classifier_challenger -- do not edit by hand.
\noindent\textbf{Run:} 2026-10-30, source \texttt{catalogue thames from PhysicalRisk-testdata}, 6 gauges scored (0 skipped), 500 storms per gauge, 5 storm folds, 6 tuning trials, 10 seeds, 20 storm bootstraps.

\begin{center}\small
\begin{tabular}{lccccccc}
\toprule
\textbf{Model} & \textbf{Brier} & \textbf{Log loss} & \textbf{AUC} & \textbf{ECE} & \textbf{Seed s.d.} & \textbf{Boot.\ s.d.} & \textbf{Mono.\ viol.} \\
\midrule
GBM (production) & 0.0204 & 0.0688 & 0.989 & 0.0078 & 0.0090 & 0.0353 & 8.3\% \\
GBM (tuned) & 0.0191 & 0.0644 & 0.990 & 0.0051 & 0.0000 & 0.0495 & 8.9\% \\
Random Forest & 0.0185 & 0.0631 & 0.990 & 0.0054 & 0.0052 & 0.0354 & 11.2\% \\
Random Forest (calib.) & 0.0199 & 0.0752 & 0.988 & 0.0061 & 0.0420 & 0.0610 & 5.3\% \\
Logistic regression & 0.0378 & 0.1438 & 0.921 & 0.0244 & 0.0000 & 0.0175 & 0.0\% \\
Flood polynomial & 0.0203 & 0.0822 & 0.979 & 0.0074 & -- & -- & 11.5\% \\
\bottomrule
\end{tabular}
\end{center}

\noindent\textbf{Decision:} best challenger by pooled Brier is Random Forest (GBM Brier 0.0204 vs 0.0185, gap 10.03\%; Wilcoxon $p = 0.016$).  Accuracy test failed, stability test failed, shape test passed: the result is \textbf{mixed} and is escalated for model-risk sign-off.

\begin{center}\small
\begin{longtable}{lrrccc}
\toprule
\textbf{Gauge} & \textbf{Storms} & \textbf{Flood rate} & \textbf{GBM Brier} & \textbf{Best chall.} & \textbf{$\Delta$ (95\% CI)} \\
\midrule\endhead
GAUGE-ef7f75ec & 500 & 12.6\% & 0.0297 & 0.0270 & +0.0027 [+0.0015, +0.0039] \\
GAUGE-066fc60e & 500 & 7.3\% & 0.0188 & 0.0166 & +0.0022 [+0.0010, +0.0034] \\
GAUGE-bc801f07 & 500 & 3.6\% & 0.0111 & 0.0102 & +0.0009 [+0.0001, +0.0018] \\
GAUGE-3324d948 & 500 & 7.3\% & 0.0210 & 0.0189 & +0.0022 [+0.0010, +0.0032] \\
GAUGE-acc1b9cc & 500 & 6.5\% & 0.0208 & 0.0197 & +0.0011 [+0.0001, +0.0022] \\
GAUGE-a03e456e & 500 & 7.1\% & 0.0208 & 0.0187 & +0.0021 [+0.0008, +0.0033] \\
\bottomrule
\end{longtable}
\end{center}
\noindent $\Delta$ is GBM minus challenger per-storm Brier; positive favours the challenger.
%

\medskip\noindent\textbf{Reading these results.}  This run is
illustrative, not sign-off evidence: it uses the first six gauges of
the 18 in the test portfolio, at 500 storms each and the default
seed and bootstrap counts.  The tuned GBM's zero seed s.d.\ does not
show better stability: tuning chose \texttt{subsample} $= 1$ for the
fold used in the stability test on all six gauges, so the fit has no
randomness for the seed to act on.  Its bootstrap s.d., which
measures sensitivity to the storms sampled, is the higher of the
GBM variants and above the Random Forest's.  Tuned parameters also
differ between folds, so its accuracy is not that of one fixed
setting.%
}{%
\noindent The comparison has not yet been run on the storm catalogue.
Until it is, the choice of GBM rests on modelling judgement and
in-sample AUC, and should be recorded as such.}


% ===================================================================
% 8. LIMITATIONS
% ===================================================================
\mkmlimitations

\begin{enumerate}[nosep]
    \item \textbf{Hydrograph-derived features only:} The four features (water level,
          hour, velocity, acceleration) are all derived from the local hydrograph.
          External features (rainfall intensity, upstream levels, antecedent
          soil moisture) could further improve discrimination but would require
          expanded data pipelines.
    \item \textbf{No spatial correlation:} Each gauge is modelled independently.
          Correlated flood events across multiple gauges are not captured.
    \item \textbf{Training noise regime:} The time-decaying noise is a modelling
          choice, not an empirical calibration. Different noise functions would
          yield different AUC profiles.
    \item \textbf{Binary classification only:} The model predicts flood/no-flood.
          Flood severity (depth, velocity, duration) is not estimated.
    \item \textbf{Stationarity assumed:} The classifier does not account for
          changing flood characteristics due to climate change or infrastructure
          modifications.
    \item \textbf{Alert-level threshold:} Training uses the FloodAlert trigger
          level. Models trained on Warning or Severe thresholds would produce
          different classifiers.
\end{enumerate}

% ===================================================================
% 9. CHANGE HISTORY
% ===================================================================
\section{Change History}

\begin{center}
\begin{tabular}{llp{7.5cm}}
\toprule
\textbf{Version} & \textbf{Date} & \textbf{Description} \\
\midrule
1.0 & 3-Mar-2026 & Initial model documentation \\
1.1 & 3-Mar-2026 & Balanced training data (1{,}000 storms), probability surface \\
1.2 & 5-Mar-2026 & Four-feature model: velocity ($\Delta w$) and acceleration
    ($\Delta^2 w$); AUC 0.9189 $\to$ 0.9557 \\
1.3 & 9-Mar-2026 & Log-transformed feature space ($\ln(w/s)$, $\ln(t/T)$,
    log-velocity, log-acceleration); AUC 0.9557 $\to$ 0.9525
    (marginal reduction, better calibration) \\
2.0 & 11-Mar-2026 & \textbf{Multi-storm upgrade:} training data replaced with
    10{,}000 StormSequence objects from MKM-SS-001; 1{,}680{,}000 vectors/gauge
    (10$\times$ v1.x); multi-peak hydrographs model antecedent-wetness dynamics;
    stressm classifiers stored in \texttt{data/output/stressm/};
    AUC 0.9525 $\to$ $\approx$0.993 (+4.2\%); scope expanded to all 52 Thames gauges \\
2.1 & 18-Mar-2026 & Stressm pipeline refactored --- classifier training now
    integrated into unified stressm pipeline with seasonal baselines from
    gaugehd monthly means \\
2.2 & 30-Oct-2026 & Added challenger comparison (GBM against Random Forest and
    others); trainer now holds out whole storms; hyperparameter table corrected
    to match the code \\
\bottomrule
\end{tabular}
\end{center}

\end{document}
